\section{Вычислительный эксперимент}\label{sec4}

Вычислительный эксперимент состоит из двух частей. Первая часть (раздел~\ref{sec4:dt}) проводится на наборе данных~\cite{Berezutskaya2022} и предназначена для оценки гемодинамической задержки $\Delta t$ и проверки корректности линейной модели, использованной при предобработке (Глава~\ref{sec3}, раздел~\ref{sec3:preproc}). Вторая часть (раздел~\ref{sec4:multimodal}) проводится на основном наборе данных~\cite{Telesford2023} и предназначена для количественной проверки предложенной двумодальной архитектуры.

\subsection{Оценка гемодинамической задержки и устойчивость линейной модели}\label{sec4:dt}

\paragraph{Набор данных и протокол.}
Используется набор данных~\cite{Berezutskaya2022}, содержащий результаты обследования 63 испытуемых, из которых для 30 доступны записи фМРТ. Все испытуемые во время регистрации смотрят один и тот же аудиовизуальный фрагмент. Параметры выборки приведены в табл.~\ref{table:dt:sample}; пример данных приведён на рис.~\ref{fig:dt:dataset}. Выборка разделена в отношении 70\%/30\% на обучающую и тестовую. На обучающей выборке вычисляется матрица параметров $\hat{\mathbf{W}}$~\eqref{eq:W_matrix}, по которой авторегрессионно прогнозируются фМРТ-сканы тестовой выборки согласно~\eqref{eq:forecast}. Метрика качества~--- $\mathrm{MSE}$ между $\hat{\mathbf{F}}_{\ell}$ и истинным $\mathbf{F}_{\ell}$, усреднённое по вокселям. Коэффициент регуляризации зафиксирован на значении $\alpha = 1000$. Значения вокселей нормализуются на отрезок $[0; 1]$ по статистикам обучающей выборки. Конфигурация вычислительной системы приведена в табл.~\ref{table:dt:pc}.

\begin{figure}[h!]
    \centering
    \includegraphics[width=0.8\textwidth]{dataset.pdf}
    \caption{Визуализация данных из~\cite{Berezutskaya2022}. Слева~--- три выбранных видеокадра, справа~--- срезы фМРТ-сканов в соответствующие моменты для трёх испытуемых. Поскольку частота кадров видео выше частоты сканирования, каждому скану соответствует группа кадров}
    \label{fig:dt:dataset}
\end{figure}

\begin{table}[h!]
\centering
\caption{Параметры выборки}\label{table:dt:sample}
\begin{tabular}{lcc}
\toprule
Параметр & Обозначение & Значение \\
\midrule
Длительность записи & $t$ & 390~с \\
Частота кадров видео & $\nu$ & 25~Гц \\
Частота сканирования фМРТ & $\mu$ & 1{,}64~Гц \\
Размеры видеокадра & $W, H, C$ & 640, 480, 3 \\
Размеры фМРТ-скана & $X, Y, Z$ & 40, 64, 64 \\
\bottomrule
\end{tabular}
\end{table}

\begin{table}[h!]
\centering
\caption{Конфигурация вычислительной системы для эксперимента по оценке $\Delta t$}\label{table:dt:pc}
\begin{tabular}{cc}
\toprule
Компонент & Описание \\
\midrule
CPU & Intel Core i7-7700 3,6~ГГц \\
GPU & NVIDIA GeForce GTX 1060 3~ГБ \\
RAM & 16~ГБ 2400~МГц \\
Накопитель & M.2 SSD \\
ОС & Windows 10 \\
\bottomrule
\end{tabular}
\end{table}

\paragraph{Пример прогноза.}
На рис.~\ref{fig:dt:example} приведены срезы истинного и предсказанного фМРТ-сканов для испытуемого~7 при $\Delta t = 5$~с. Поскольку значения вокселей нормализованы на $[0; 1]$, ошибка порядка $10^{-3}$ соответствует достаточно точному прогнозу.

\begin{figure}[h!]
    \centering
    \subfloat[Истинный скан]{\includegraphics[width=0.32\textwidth]{sub-07-5-1-1000-37-20-_-_-test.png}}
    \hfill
    \subfloat[Прогноз]{\includegraphics[width=0.32\textwidth]{sub-07-5-1-1000-37-20-_-_-predicted.png}}
    \hfill
    \subfloat[Разность]{\includegraphics[width=0.32\textwidth]{sub-07-5-1-1000-37-20-_-_-difference.png}}
    \caption{Срезы фМРТ-сканов из тестовой выборки. Истинный и прогнозируемый сканы, а также их попиксельная разность. Ошибка порядка $10^{-3}$ свидетельствует о точности прогноза}
    \label{fig:dt:example}
\end{figure}

\paragraph{Зависимость качества прогноза от $\Delta t$.}
Исследована зависимость ошибки $\mathrm{MSE}$ от гиперпараметра $\Delta t$ на испытуемом~47. Левый график рис.~\ref{fig:dt:mse-dt} показывает зависимость, вычисленную на всех вокселях скана: чёткий минимум не наблюдается, поскольку доля фоновых и неактивированных вокселей вносит существенный шум. Для устранения этой шумовой составляющей локализуется затылочная доля~--- область, ответственная за обработку визуальной информации~\cite{anderson2006}; из локализованного объёма выбираются 3\% наиболее изменчивых вокселей (отсортированных по убыванию суммарного абсолютного изменения значений~\eqref{eq:mask_score}). Правый график рис.~\ref{fig:dt:mse-dt} демонстрирует чёткий минимум при $\Delta t \approx 5$~с, что согласуется с известным диапазоном гемодинамической задержки BOLD-сигнала~\cite{Bandettini1992}. Локализация затылочной доли проиллюстрирована на рис.~\ref{fig:dt:local}. Полученное значение $\Delta t = 5$~с фиксируется и используется в основном эксперименте при формировании триплетов согласно~\eqref{eq:shift}.

\begin{figure}[h!]
    \centering
    \subfloat[Истинный скан]{\includegraphics[width=0.32\textwidth]{sub-47-5-1-1000-37-20-_-_-test.png}}
    \hfill
    \subfloat[Прогноз]{\includegraphics[width=0.32\textwidth]{sub-47-5-1-1000-37-20-_-_-predicted.png}}
    \hfill
    \subfloat[Разность]{\includegraphics[width=0.32\textwidth]{sub-47-5-1-1000-37-20-_-_-difference.png}}
    \caption{Локализация наиболее активной зоны при просмотре видеостимула. Выделенная область~--- затылочная доля, что подтверждает корректность локализации}
    \label{fig:dt:local}
\end{figure}

\begin{figure}[h!]
    \centering
    \includegraphics[width=\textwidth]{mse_dt.pdf}
    \caption{Зависимость $\mathrm{MSE}$ от гемодинамической задержки $\Delta t$. Слева~--- расчёт на всех вокселях скана: чёткий минимум отсутствует из-за вклада фоновых вокселей. Справа~--- расчёт на 3\% наиболее изменчивых вокселей затылочной доли: наблюдается чёткий минимум при $\Delta t \approx 5$~с}
    \label{fig:dt:mse-dt}
\end{figure}

\paragraph{Распределение весов линейной модели.}
Для проверки осмысленности обученной модели исследовано распределение компонентов вектора весов $\hat{\mathbf{w}}_{ijk}$ из~\eqref{eq:ridge_sol}, усреднённое по всем вокселям испытуемого~4 (рис.~\ref{fig:dt:weights}). Распределение не сосредоточено около одного значения, что свидетельствует о невырожденности модели и согласуется с интерпретацией: при просмотре видеостимула внимание испытуемого направлено на различные элементы кадра, и линейная модель отражает это в неоднородной структуре весов по компонентам визуального эмбеддинга.

\begin{figure}[h!]
    \centering
    \includegraphics[width=0.65\textwidth]{distribution.pdf}
    \caption{Распределение компонент вектора весов линейной модели. Невырожденность распределения свидетельствует о том, что модель захватывает разнообразные визуальные признаки}
    \label{fig:dt:weights}
\end{figure}

\paragraph{Инвариантность весов относительно испытуемых.}
Проверяется гипотеза о том, что веса линейной модели~\eqref{eq:voxel_linear} инвариантны относительно испытуемых: матрица $\hat{\mathbf{W}}$, обученная на одном испытуемом, может использоваться для прогнозирования фМРТ другого без существенной потери качества. Для пары испытуемых~4 и~7 значения $\mathrm{MSE}$ при использовании <<своей>> и <<чужой>> матрицы весов практически совпадают (табл.~\ref{table:dt:inv}).

\begin{table}[h!]
\centering
\caption{Проверка инвариантности весов линейной модели для пары испытуемых~4 и~7}\label{table:dt:inv}
\begin{tabular}{lccc}
\toprule
Матрица весов & Своя & Чужая (от другого испытуемого) & Разность \\
\midrule
$\mathrm{MSE}$ & $6{,}3912 \cdot 10^{-5}$ & $6{,}3911 \cdot 10^{-5}$ & $1{,}37 \cdot 10^{-9}$ \\
\bottomrule
\end{tabular}
\end{table}

Эксперимент проведён для всех пар испытуемых; результаты сведены в тепловую карту (рис.~\ref{fig:dt:heatmap}): для каждой пары вычислена $\mathrm{MAPE}$ между <<своим>> $\mathrm{MSE}$ (с собственной матрицей весов) и <<чужим>> $\mathrm{MSE}$ (с матрицей весов другого испытуемого). Отклонения порядка 1\% подтверждают гипотезу: линейная модель захватывает универсальную часть связи между видеостимулом и BOLD-откликом, общую для разных испытуемых. Этот результат поддерживает обоснованность фиксации единого значения $\Delta t$ при предобработке триплетов основного эксперимента.

\begin{figure}[h!]
    \centering
    \includegraphics[width=0.5\textwidth]{heatmap.pdf}
    \caption{Тепловая карта $\mathrm{MAPE}$ при подмене матрицы весов. Каждая строка и каждый столбец соответствуют испытуемому. Положительные значения означают, что использование чужой матрицы увеличивает $\mathrm{MSE}$ относительно своей}
    \label{fig:dt:heatmap}
\end{figure}

\paragraph{Контроль на неинформативных признаках.}
Для подтверждения того, что модель действительно использует визуальную информацию из видеокадров, а не подгоняется к структуре приращений фМРТ, проведён следующий эксперимент. Матрица визуальных эмбеддингов $\boldsymbol{\Psi}$ заменяется на матрицу из единиц (неинформативные признаки), после чего модель переобучается по~\eqref{eq:ridge_sol} и используется для авторегрессионного прогноза~\eqref{eq:forecast} последнего фМРТ-скана тестовой выборки путём суммирования прогнозируемых приращений начиная с первого скана. На рис.~\ref{fig:dt:recover} приведены результаты для исходных видеопризнаков, на рис.~\ref{fig:dt:random}~--- для неинформативных. Численные результаты (испытуемый~35) представлены в табл.~\ref{table:dt:random}: $\mathrm{MSE}$ на неинформативных признаках почти в четыре раза больше, что подтверждает наличие значимой корреляции между видеостимулом и BOLD-сигналом.

\begin{figure}[h!]
    \centering
    \subfloat[Истинный скан]{\includegraphics[width=0.32\textwidth]{sub-35-5-1-1000--1-20-_-_-recovered-test.png}}
    \hfill
    \subfloat[Прогноз]{\includegraphics[width=0.32\textwidth]{sub-35-5-1-1000--1-20-_-_-recovered-predicted.png}}
    \hfill
    \subfloat[Разность]{\includegraphics[width=0.32\textwidth]{sub-35-5-1-1000--1-20-_-_-recovered-difference.png}}
    \caption{Срезы фМРТ-сканов из тестовой выборки при использовании оригинальных видеопризнаков}
    \label{fig:dt:recover}
\end{figure}

\begin{figure}[h!]
    \centering
    \subfloat[Истинный скан]{\includegraphics[width=0.32\textwidth]{noised_sub-35-5-1-1000--1-20-_-_-recovered-test.png}}
    \hfill
    \subfloat[Прогноз]{\includegraphics[width=0.32\textwidth]{noised_sub-35-5-1-1000--1-20-_-_-recovered-predicted.png}}
    \hfill
    \subfloat[Разность]{\includegraphics[width=0.32\textwidth]{noised_sub-35-5-1-1000--1-20-_-_-recovered-difference.png}}
    \caption{Срезы фМРТ-сканов при использовании неинформативных признаков. $\mathrm{MSE}$ существенно выше, что подтверждает значимость видеопризнаков}
    \label{fig:dt:random}
\end{figure}

\begin{table}[h!]
\centering
\caption{Сравнение качества прогноза при использовании оригинальных и неинформативных признаков}\label{table:dt:random}
\begin{tabular}{cccc}
\toprule
Признаки & Оригинальные & Неинформативные & Разность \\
\midrule
$\mathrm{MSE}$ & $3{,}67 \cdot 10^{-4}$ & $1{,}39 \cdot 10^{-3}$ & $1{,}02 \cdot 10^{-3}$ \\
\bottomrule
\end{tabular}
\end{table}

\paragraph{Выводы по эксперименту с линейной моделью.}
Эксперименты на наборе данных~\cite{Berezutskaya2022} подтверждают: значимая корреляция между видеостимулом и BOLD-сигналом (качество прогноза на оригинальных видеопризнаках существенно превышает качество на неинформативных); оценка гемодинамической задержки $\Delta t \approx 5$~с, согласующаяся с известным медицинским диапазоном~\cite{Bandettini1992}; инвариантность весов линейной модели относительно испытуемых, что обосновывает использование единого значения $\Delta t$ для всех испытуемых при формировании триплетов основного эксперимента; невырожденное распределение весов модели, согласующееся с содержательной интерпретацией. Полученное значение $\Delta t = 5$~с фиксируется при предобработке триплетов основного эксперимента (раздел~\ref{sec4:multimodal}).

\subsection{Мультимодальная реконструкция изображений}\label{sec4:multimodal}

\paragraph{Набор данных.}
В качестве основного набора данных используется открытый набор с одновременной регистрацией фМРТ и ЭЭГ при предъявлении визуальных стимулов~\cite{Telesford2023}. В набор включены 22 здоровых испытуемых в возрасте 23--51 года; каждый испытуемый прошёл два сканирования. Стимулы охватывают состояние покоя, мерцающее шахматное поле, абстрактную анимацию Inscapes (компьютерно сгенерированную последовательность трёхмерных фигур с плавными переходами), а также короткие видеоролики: фрагменты мультфильмов \textit{The Present} и \textit{Despicable Me} и фрагменты документальных фильмов с обезьянами. Натуралистический характер стимулов способствует устойчивому межсубъектному нейронному отклику. Синхронизация модальностей обеспечивается записью триггеров сканера в поток ЭЭГ. Предобработка модальностей выполняется согласно процедуре, описанной в Главе~\ref{sec3} (раздел~\ref{sec3:preproc}). Параметры регистрации приведены в табл.~\ref{table:m:params}. Пример индивидуальной маски мозга, построенной алгоритмом Оцу, приведён на рис.~\ref{fig:m:masking}. Триплеты $(\mathbf{F}_i, \mathbf{E}_i, \mathbf{I}_i)$ формируются с учётом фиксированной гемодинамической задержки $\Delta t = 5$~с согласно~\eqref{eq:shift}, обоснованной экспериментом раздела~\ref{sec4:dt}.

\begin{table}[h!]
    \centering
    \caption{Параметры регистрации модальностей в наборе данных~\cite{Telesford2023}}
    \label{table:m:params}
    \resizebox{\textwidth}{!}{\begin{tabular}{llll}
    \toprule
    Модальность & Оборудование & Частота дискретизации & Каналы \\
    \midrule
    ЭЭГ & Brain Products BrainCap MR & 5000~Гц & 61 кортикальный электрод \\
    фМРТ & Siemens 3.0~T TIM Trio & 0,476~Гц & 12-канальная головная катушка \\
    \bottomrule
    \end{tabular}}
\end{table}

\begin{figure}[h!]
    \centering
    \includegraphics[width=0.8\textwidth]{figure_4.png}
    \caption{Индивидуальная бинарная маска мозга для одного из испытуемых. Включённые в маску воксели выделены красным цветом. Маска формируется на основе наиболее изменчивых вокселей во времени и удаляет фоновые воксели, не несущие информации о визуальном стимуле}
    \label{fig:m:masking}
\end{figure}

\paragraph{Параметры энкодеров.}
\emph{Энкодер фМРТ:} индивидуальная линейная проекция отображает векторизованные воксели, отобранные бинарной маской испытуемого, в латентное пространство размерности 4096; далее остаточный MLP с тремя полносвязными слоями последовательно $4096 \to 2048 \to 2048 \to 1024$, нелинейностью GELU, нормализацией по слою и dropout с вероятностью 0{,}1; размерность выходного эмбеддинга $d_f = 1024$. \emph{Энкодер ЭЭГ:} вход~--- тензор размера $61 \times 525$ (60 каналов корковых электродов и один канал индивидуального обучаемого эмбеддинга испытуемого); три свёртки (временная с ядром $(1, 32)$ и 256 фильтрами, пространственная с ядром $(61, 1)$ и 512 фильтрами, глобальная с ядром $(1, 525)$ и 1024 фильтрами); MLP-голова $1024 \to 2048 \to 1024$; размерность выходного эмбеддинга $d_e = 1024$. \emph{Модуль объединения:} двухслойный MLP, конкатенированное представление размерности $2048$ расширяется до $4096$ и сжимается до целевой размерности $d = 1024$, совпадающей с размерностью CLIP-пространства.

\paragraph{Качественное сравнение реконструкций.}
Качественное сравнение реконструкций приведено на рис.~\ref{fig:m:reconstructions}. Строки соответствуют различным режимам: (1)~исходный визуальный стимул; (2)~реконструкция из CLIP-эмбеддинга исходного изображения через предобученную SDXL, представляющая верхнюю границу качества, достижимую в выбранной схеме генерации; (3)~реконструкция предложенным методом без априорной диффузионной модели; (4)~реконструкция предложенным методом с априорной диффузионной моделью. Включение априорной диффузионной модели визуально повышает семантическое соответствие реконструкции исходному стимулу.

\begin{figure}[h!]
    \centering
    \includegraphics[width=1.0\textwidth]{figure_5.pdf}
    \caption{Качественное сравнение реконструкций изображений по сигналам фМРТ-ЭЭГ. Строки: (1)~исходный стимул; (2)~предельное качество (SDXL по CLIP-эмбеддингу оригинала); (3)~предложенный метод без априорной диффузионной модели; (4)~предложенный метод с априорной диффузионной моделью}
    \label{fig:m:reconstructions}
\end{figure}

\paragraph{Количественное сравнение и ablation.}
Для количественной оценки используется метрика CLIP-Score~\eqref{eq:clip_score}. Метрики вычислялись на тестовой выборке из видеостимулов \texttt{monkey1}, \texttt{monkey2}, \texttt{monkey5} с усреднением по 22 испытуемым и по объектам. Сравниваются шесть конфигураций: (1)~\textbf{ЭЭГ}~--- только энкодер $\mathcal{E}_e$, модуль объединения отключен, $\mathbf{Z}_e$ напрямую используется как условие для генерации, априорная модель не применяется; (2)~\textbf{фМРТ}~--- аналогично с энкодером $\mathcal{E}_f$; (3)~\textbf{фМРТ-ЭЭГ}~--- оба энкодера и модуль объединения, без априорной модели; (4)~\textbf{ЭЭГ \& Diffusion Prior}~--- энкодер ЭЭГ с априорной диффузионной моделью; (5)~\textbf{фМРТ \& Diffusion Prior}~--- энкодер фМРТ с априорной диффузионной моделью; (6)~\textbf{фМРТ-ЭЭГ \& Diffusion Prior}~--- полная архитектура. Результаты сведены в табл.~\ref{tab:m:clip}.

\begin{table}[h!]
    \centering
    \caption{Метрика CLIP-Score для шести конфигураций модели на трёх видеостимулах. Усреднение по 22 испытуемым; нижним индексом указаны стандартные отклонения. Наилучший результат в столбце выделен жирным шрифтом}
    \label{tab:m:clip}
    \begin{tabular}{lccc}
        \toprule
        & \multicolumn{3}{c}{Видеостимул} \\
        \cline{2-4}
        Конфигурация & \texttt{monkey1} & \texttt{monkey2} & \texttt{monkey5} \\
        \midrule
        ЭЭГ                       & $0{,}231_{\pm 0{,}004}$ & $0{,}262_{\pm 0{,}004}$ & $0{,}224_{\pm 0{,}005}$ \\
        фМРТ                      & $0{,}254_{\pm 0{,}004}$ & $0{,}270_{\pm 0{,}003}$ & $0{,}242_{\pm 0{,}003}$ \\
        фМРТ-ЭЭГ                  & $0{,}310_{\pm 0{,}003}$ & $0{,}312_{\pm 0{,}004}$ & $0{,}284_{\pm 0{,}003}$ \\
        ЭЭГ \& Diffusion Prior    & $0{,}542_{\pm 0{,}005}$ & $0{,}557_{\pm 0{,}003}$ & $0{,}538_{\pm 0{,}003}$ \\
        фМРТ \& Diffusion Prior   & $0{,}564_{\pm 0{,}005}$ & $0{,}571_{\pm 0{,}004}$ & $0{,}551_{\pm 0{,}005}$ \\
        фМРТ-ЭЭГ \& Diffusion Prior & $\mathbf{0{,}668}_{\pm 0{,}003}$ & $\mathbf{0{,}647}_{\pm 0{,}005}$ & $\mathbf{0{,}644}_{\pm 0{,}003}$ \\
        \bottomrule
    \end{tabular}
\end{table}

Двумодальная архитектура (фМРТ-ЭЭГ) систематически превосходит обе одномодальные конфигурации на всех трёх видеостимулах, что подтверждает гипотезу о комплементарности пространственной информации фМРТ и временной информации ЭЭГ. Включение априорной диффузионной модели повышает CLIP-Score для всех конфигураций приблизительно вдвое; этот эффект согласуется с наблюдениями в одномодальной литературе~\cite{scotti2023reconstructingmindseyefmritoimage}: априорная модель устраняет разрыв между распределением контрастивно обученных нейронных эмбеддингов и распределением CLIP-эмбеддингов реальных изображений. Наилучшее качество достигается в полной конфигурации (фМРТ-ЭЭГ \& Diffusion Prior).

\paragraph{Достоверность и ограничения.}
Достоверность представленных результатов обеспечивается следующими факторами: использование открытого набора данных с протоколом, прошедшим этическую экспертизу; усреднение по 22 испытуемым с указанием стандартных отклонений; систематическое ablation-сравнение шести конфигураций; сравнение с верхней границей качества, достижимой при условии на CLIP-эмбеддинг исходного изображения (рис.~\ref{fig:m:reconstructions}, строка~2); независимая проверка корректности значения $\Delta t$ на отдельном наборе данных (раздел~\ref{sec4:dt}). Использование CLIP-пространства в качестве промежуточного представления ограничивает реконструкцию семантическим уровнем: точное восстановление пиксельных деталей, текстур или низкоуровневых визуальных характеристик не гарантируется. Экспериментальная проверка проведена на одном датасете с видеостимулами определённого типа; перенос результатов на статичные изображения, абстрактные паттерны или иные условия регистрации требует дополнительной валидации. Метод обучен и протестирован на данных здоровых испытуемых; применение к пациентам с патологиями зрительной системы или нарушениями когнитивных функций выходит за рамки настоящей работы.
